%% ==========================================================================
%%  Approach 3 --- the replica transform and the peel reformulation
%%
%%  Source: update_2/PS1_note_replica_peel_transform.md (v1, 2026-08-02).
%%  The replica construction is the author's; the downstream development was
%%  worked out in the session of that date.
%%
%%  STATUS.  Theorem H and its corollaries are proved.  The peel dead-end
%%  computation is machine-verified (update_2/peel3.py reproduces every number
%%  in this section exactly).  Conjecture H1 is REFUTED; Conjecture H1' is
%%  OPEN and is the statement that would imply Conjecture 2.
%%
%%  This section is deliberately less polished than the rest of the parked
%%  material.  Sub-sections marked [Conjecture] or [Observation] are not
%%  proved and must be labelled as such in anything shown to a reader.
%%
%%  Terms defined here are flagged for promotion into
%%  ../glossary_fair_division_subsidies.md, which has NOT yet been updated:
%%  replica instance, type map, dual valuation, coverage allocation,
%%  duplicate copy, peel process, workload profile, owner-candidate set,
%%  survivor, spectator-free peel, conditioned-remainder principle, deadlock.
%% ==========================================================================
\section{Approach 3: The Replica Transform and the Peel Reformulation}
\label{sec:approach3}

Both preceding approaches fail because of \emph{when} decisions are made.
Item-by-item insertion (\Cref{sec:approach1}) commits a chore to an owner and
cannot take it back; selecting a utilitarian optimum (\Cref{sec:approach2})
commits to the wrong allocation outright. This section changes coordinates so
that the commitment is deferred. The chore problem is re-expressed as a
\emph{goods} problem on a replicated item set, and then read dynamically as a
process that relieves agents of chores one at a time. Ownership of a chore is
decided only by the last move that touches it.

The payoff is that the proof orientation of~\cite{BKNS22} becomes correct again
and the obstruction of Theorem~\ref{thm:obstruction} ceases to bind. The price
is a new constraint --- coverage --- which is then shown to be non-vacuous.

% --------------------------------------------------------------------------
\subsection{The replica instance}
\label{sec:replica}

Fix a negative dichotomous instance with cost functions $\cost_i$ on chores
$\items = \set{a_1,\dots,a_m}$. Build a \emph{replica instance} $\hat{I}$ whose
item set carries $n-1$ copies of a good $b_j$ for each chore $a_j$:
\[
  \hat{\items} \;:=\; \set{\, b_j^{(1)}, \dots, b_j^{(n-1)} \;:\; j \in \items \,}.
\]
For $S \subseteq \hat{\items}$ let $\tau(S) \subseteq \items$ be the set of
\emph{types} occurring in $S$, and define
\[
  \hat{v}_i(S) \;:=\; \cost_i(\items) \;-\; \cost_i\big(\items \setminus \tau(S)\big).
\]

The reading is that \emph{handing agent $i$ a copy of $b_j$ relieves $i$ of
chore $a_j$}. Each chore ends with exactly one owner, so exactly $n-1$ agents
must be relieved of it --- hence $n-1$ copies --- and no agent may be relieved
of the same chore twice, so no agent may hold two copies of one type. That last
clause is load-bearing rather than cosmetic: it is what makes $\hat{v}_i$
well defined as a dual, since a second copy of a type has marginal $0$ and never
$+1$ again.

\begin{remark}[The naive reading is false]
\label{rem:naive}
If one instead declares a copy of $b_j$ to be worth $+1$ to whoever receives it,
the construction collapses. An agent indifferent to $a_j$ would gain from
$b_j$, every $\hat{v}_i$ would reduce to $\abs{\tau(S)}$, and the induced envy
graph would be $\hat{\edgew{}}(i,k) = \abs{B_i} - \abs{B_k}$ --- the pure
cardinality term already isolated in Theorem~\ref{thm:sizeshift}, carrying none
of the cost information. The correct reading is \emph{marginal}: a copy of $b_j$
is worth to agent $i$ exactly $i$'s marginal cost of $a_j$ on what remains of
$i$'s workload.
\end{remark}

This also completes Remark~\ref{rem:complement}. There the dual
$\tilde{v}_i(S) = \cost_i(\items) - \cost_i(\items \setminus S)$ settled $n = 2$
and broke at $n \ge 3$ because $\items \setminus \bundle{i}$ stops being a
bundle of the partition. Replication is precisely the repair: with $n-1$ copies
of every type, the $n$ complements $\items \setminus \bundle{1}, \dots,
\items \setminus \bundle{n}$ become \emph{simultaneously} realisable as disjoint
bundles.

% --------------------------------------------------------------------------
\subsection{The transform}
\label{sec:thmH}

Call an allocation $B = (B_1,\dots,B_n)$ of all of $\hat{\items}$ a
\emph{coverage allocation} if no $B_i$ contains two copies of the same type.

\begin{theorem}[Replica transform]
\label{thm:replica}
Let $\cost_1,\dots,\cost_n$ be monotone with $\cost_i(\emptyset) = 0$ and all
marginals in $\set{0,1}$. Then
\begin{enumerate}
  \item[(i)] each $\hat{v}_i$ is a dichotomous valuation on $\hat{\items}$ in
    the sense of~\cite{BKNS22};
  \item[(ii)] $\Phi : \alloc \mapsto B$, where $B_i$ receives one copy of each
    type in $\items \setminus \bundle{i}$, is a bijection between partitions of
    $\items$ and coverage allocations of $\hat{\items}$, up to relabelling
    copies of a type;
  \item[(iii)] for corresponding $\alloc$ and $B$ the envy graphs agree arc for
    arc: $\hat{\edgew{}}_B(i,k) = \edgew{\alloc}(i,k)$ for all $i,k$.
\end{enumerate}
\end{theorem}

\begin{proof}
\emph{(i)} Normalisation: $\hat{v}_i(\emptyset) = \cost_i(\items) -
\cost_i(\items) = 0$. Monotonicity: $S \subseteq S'$ gives $\tau(S) \subseteq
\tau(S')$, hence $\items \setminus \tau(S) \supseteq \items \setminus \tau(S')$,
and $\cost_i$ is monotone, so $\hat{v}_i(S) \le \hat{v}_i(S')$. Marginals: take
$b \notin S$ of type $j$. If $j \in \tau(S)$ then $\tau(S \cup \set{b}) =
\tau(S)$ and the marginal is $0$. Otherwise, writing $T := \items \setminus
\tau(S)$, we have $\items \setminus \tau(S \cup \set{b}) = T \setminus \set{j}$
with $j \in T$, so
\[
  \hat{v}_i(S \cup \set{b}) - \hat{v}_i(S)
  \;=\; \cost_i(T) - \cost_i(T \setminus \set{j}) \;\in\; \set{0,1},
\]
because $\cost_i$ has marginals in $\set{0,1}$.

\emph{(ii)} Forward: type $j$ is absent from exactly one bundle of $\alloc$, its
owner's, so exactly $n-1$ agents require a copy of $b_j$ and the $n-1$ available
copies are routed injectively. Which physical copy goes where is immaterial,
since $\hat{v}_i$ depends on $S$ only through $\tau(S)$. By construction no
$B_i$ holds a duplicate and $\bigcup_i B_i = \hat{\items}$. Backward: in a
coverage allocation the $n-1$ copies of type $j$ lie in $n-1$ \emph{distinct}
bundles, so exactly one agent's bundle contains no copy of $j$; setting
$\Psi(B)_i := \items \setminus \tau(B_i)$ therefore puts $j$ in exactly one
$\Psi(B)_i$, making $\Psi(B)$ a partition. Finally $\Psi(\Phi(\alloc))_i =
\items \setminus (\items \setminus \bundle{i}) = \bundle{i}$, and
$\Phi(\Psi(B))$ agrees with $B$ up to which copy of a type sits where.

\emph{(iii)} For corresponding $\alloc, B$ we have $\tau(B_k) = \items \setminus
\bundle{k}$, so $\hat{v}_i(B_k) = \cost_i(\items) - \cost_i(\bundle{k})$ and
\[
  \hat{\edgew{}}_B(i,k) = \hat{v}_i(B_k) - \hat{v}_i(B_i)
  = \big(\cost_i(\items) - \cost_i(\bundle{k})\big)
  - \big(\cost_i(\items) - \cost_i(\bundle{i})\big)
  = \cost_i(\bundle{i}) - \cost_i(\bundle{k}) = \edgew{\alloc}(i,k). \qedhere
\]
\end{proof}

Part~(i) is worth stating on its own: \emph{the dual of a dichotomous cost
function is a dichotomous value function}, so the class is closed under this
duality and the hypotheses of~\cite{BKNS22} apply to $\hat{I}$ verbatim, with no
additivity, submodularity or subadditivity assumed anywhere. Because the graphs
of part~(iii) are equal, envy-freeability, every $\pathw{}(i)$, the whole set of
feasible subsidy vectors and the $\set{0,1}$ target all transfer in both
directions without loss.

\begin{corollary}[Coverage is the entire gap]
\label{cor:coverage}
Conjecture~\ref{conj:main} holds if and only if every replica instance
$\hat{I}$ admits a \emph{coverage} allocation $B$ with
$\pathw{B}(i) \le 1$ for all $i$.
\end{corollary}

Applying~\cite{BKNS22} to $\hat{I}$ as a black box is legal by
Theorem~\ref{thm:replica}(i) and returns some $B$ with $\hat{\subsidy} \in
\set{0,1}^n$ and total at most $n-1$. But it is free to place two copies of one
type in one bundle. If it does, that type is absent from at least two bundles,
the chore has at least two owners, and $\Psi(B)$ is not a partition. Coverage is
the sole residual difficulty --- and, by \Cref{sec:deadends}, a real one.

\begin{lemma}[Duplicate extraction is envy-neutral]
\label{lem:dupfree}
Removing from $B_x$ a copy whose type occurs at least twice in $B_x$ leaves
$\tau(B_x)$ unchanged, hence leaves every $\hat{v}_k(B_x)$ unchanged, hence
leaves the entire envy graph unchanged.
\end{lemma}

\begin{proof}
Immediate, since $\hat{v}_k(S)$ depends on $S$ only through $\tau(S)$.
\end{proof}

So any output of~\cite{BKNS22} may be stripped to a duplicate-free partial
allocation for free. The difficulty is never in removing duplicates; it is in
\emph{re-inserting} the stripped copies into agents that do not already hold
that type.

% --------------------------------------------------------------------------
\subsection{The peel process}
\label{sec:peel}

Read a coverage allocation as being built one copy at a time and translate back
into chores.

\begin{definition}[Peel process]
\label{def:peel}
A \emph{state} is a \emph{workload profile} $\workprofile = (\workprofile_1,
\dots, \workprofile_n)$ with $\workprofile_i \subseteq \items$, such that the
\emph{owner-candidate set} $S_j := \set{i : j \in \workprofile_i}$ is nonempty
for every $j \in \items$. Read $\workprofile_i$ as ``agent $i$ is still on the
hook for $\workprofile_i$''.
\begin{itemize}
  \item \textbf{Root:} $\workprofile_i = \items$ for every $i$.
  \item \textbf{Peel} $\mathrm{peel}(x,j)$: legal iff $j \in \workprofile_x$ and
    $\abs{S_j} \ge 2$; it sets $\workprofile_x \leftarrow \workprofile_x
    \setminus \set{j}$.
  \item \textbf{Permutation:} replace $\workprofile$ by $(\workprofile_{\sigma(1)},
    \dots, \workprofile_{\sigma(n)})$.
  \item \textbf{Terminal:} $\abs{S_j} = 1$ for every $j$; the unique element of
    $S_j$ is the \emph{survivor}, i.e.\ the owner of $j$, and the induced
    allocation is $\bundle{i} = \set{j : S_j = \set{i}}$.
  \item \textbf{Graph:} $\edgew{\workprofile}(i,k) = \cost_i(\workprofile_i) -
    \cost_i(\workprofile_k)$. The root graph is identically zero.
\end{itemize}
A state is \emph{legal} if its graph has no positive-weight cycle and
$\pathw{\workprofile}(i) \le 1$ for every $i$.
\end{definition}

By Theorem~\ref{thm:replica} a peel state is exactly a duplicate-free partial
allocation of $\hat{\items}$, a peel is exactly the insertion of one good into
one bundle, a permutation is exactly a bundle reassignment, and terminal states
are exactly coverage allocations. Nothing is added or lost in the translation.

\begin{lemma}[Peel dynamics]
\label{lem:peel}
Write $\mu_k := \cost_k(\workprofile_x) - \cost_k(\workprofile_x \setminus
\set{j})$ for $k$'s marginal cost of $j$ on $x$'s residual workload. Then
$\mathrm{peel}(x,j)$ changes only the arcs incident to $x$, by
\[
  \Delta \edgew{}(k,x) \;=\; +\mu_k \;\in\; \set{0,1}
  \qquad\text{and}\qquad
  \Delta \edgew{}(x,k) \;=\; -\mu_x \;\in\; \set{-1,0}
  \qquad (k \ne x).
\]
\end{lemma}

\begin{proof}
Only $\workprofile_x$ changes. For $k \ne x$, $\edgew{}(k,x) =
\cost_k(\workprofile_k) - \cost_k(\workprofile_x)$ and $\cost_k(\workprofile_x)$
drops by $\mu_k$, so the arc rises by $\mu_k$. For arcs out of $x$,
$\edgew{}(x,k) = \cost_x(\workprofile_x) - \cost_x(\workprofile_k)$ and
$\cost_x(\workprofile_x)$ drops by $\mu_x$. Arcs between two agents other than
$x$ involve neither $\workprofile_x$ nor $\cost_x$.
\end{proof}

\paragraph{The orientation reversal.}
This is the point of the whole construction. Compare
Lemma~\ref{lem:peel} against Lemma~\ref{lem:arcupdate}:

\begin{center}
\begin{tabular}{@{}llll@{}}
\toprule
frame & move & in-arcs at $x$ & $x$ should be \\
\midrule
goods~\cite{BKNS22}      & insert a good   & rise          & maximally subsidised \\
chores, insertion (\S\ref{sec:approach1}) & insert a chore  & fall          & unsubsidised, $\subsidy_x = 0$ \\
chores, \textbf{peel}    & relieve of a chore & \textbf{rise} & maximally subsidised \\
\bottomrule
\end{tabular}
\end{center}

The obstruction of Theorem~\ref{thm:obstruction} was exactly this sign clash.
The repair subroutine of~\cite{BKNS22} navigates to an agent whose tentative
subsidy reached $2$; in the goods world that agent is automatically maximally
subsidised, which is the set the argument needs, whereas in the chore-insertion
world she automatically has $\subsidy_u = 1$ and is precisely the agent who must
\emph{not} receive the chore. In the peel frame the quantity handed out is
\emph{relief}, relief correctly flows to the most envious agents, and the
orientation points the right way again.

\begin{lemma}[Spectator-free peels are safe]
\label{lem:specfree}
If $\mu_k = 0$ for every $k \ne x$, then $\mathrm{peel}(x,j)$ weakly decreases
every arc weight, hence $\pathw{}$ pointwise, and preserves envy-freeability ---
whatever $\mu_x$ may be.
\end{lemma}

\begin{proof}
By Lemma~\ref{lem:peel} the arcs into $x$ are unchanged, the arcs out of $x$
change by $-\mu_x \le 0$, and all other arcs are untouched. Path and cycle
weights are sums of arc weights.
\end{proof}

Lemma~\ref{lem:specfree} is the peel-frame replacement for
Theorem~\ref{thm:freeinsert}, and note where the safety condition has moved. In
the insertion frame a chore that was free \emph{for the receiver} was
unconditionally safe; here what matters is that the chore be free \emph{for
everyone else}. That is the goods-side safety structure of~\cite{BKNS22}, as the
identification predicts. The hard case is now a pair $(x,j)$ for which some
other agent $k$ has $\mu_k = 1$.

What does \emph{not} transfer is the coverage constraint, which has no analogue
in~\cite{BKNS22}: any good may go in any bundle there, whereas a peel demands
$x \in S_j$ and a type is peelable only while $\abs{S_j} \ge 2$.

% --------------------------------------------------------------------------
\subsection{The insertion obstruction dissolves}
\label{sec:dissolve}

Theorem~\ref{thm:obstruction} remains true --- it is a theorem about the
insertion template, and the peel frame is a different template, not a repair of
that one. What follows is that its witness no longer obstructs.

Recall the witness: $\items = \set{a_1,a_2,g}$ with $\cost_1(S) =
\max(0,\abs{S}-1)$ and $\cost_2 = \cost_3 = \abs{S}$.

\paragraph{The trap state is refused before it can form.}
The peel-frame counterpart of the trap is $\workprofile =
(\set{a_1,a_2,g}, \set{g}, \set{g})$, and
\[
  \edgew{\workprofile}(1,2) \;=\; \cost_1(\set{a_1,a_2,g}) - \cost_1(\set{g})
  \;=\; 2 - 0 \;=\; 2 ,
\]
so $\pathw{\workprofile}(1) = 2$ and the state is already illegal. No legal peel
schedule passes through it.

\begin{remark}[The conditioned-remainder principle]
\label{rem:conditioned}
The two frames disagree because they compare different things. At a stage where
the finished types form a partial assignment $\alloc^{(t)}$ and the untouched
types form $R$, every agent is still on the hook for $R$, so the peel-frame arc
is
\[
  \cost_i\big(\alloc^{(t)}_i \cup R\big) - \cost_i\big(\alloc^{(t)}_k \cup R\big),
  \qquad\text{not}\qquad
  \cost_i\big(\alloc^{(t)}_i\big) - \cost_i\big(\alloc^{(t)}_k\big).
\]
Every comparison is conditioned on the common unfinished remainder. For additive
costs the $R$ terms cancel and the two frames coincide --- consistent with the
additive case already being closed by Theorem~\ref{thm:binadd}. For non-additive
costs they differ, and on the witness it is exactly the supermodularity of
$\cost_1$ that the conditioning catches: with $g$ still pending, loading agent~1
with both of $a_1,a_2$ already costs $2$, and the peel invariant sees it. \emph{The
insertion invariant was blind to pending load; the peel invariant has hindsight
built in.} This is the substantive content of the change of coordinates.
\end{remark}

\paragraph{A legal schedule reaches the zero-subsidy allocation.}
Six peels, every one of them the hard case ($\mu_k = 1$ for all $k$ --- this
instance never offers a spectator-free move), with the invariant holding
throughout. Verified by \texttt{peel.py}.

\begin{center}
\begin{tabular}{@{}clccc@{}}
\toprule
step & move & resulting $\workprofile$ & $\pathw{}$ & $\sum_i \pathw{}(i)$ \\
\midrule
0 & root & $(a_1a_2g,\; a_1a_2g,\; a_1a_2g)$ & $(0,0,0)$ & 0 \\
1 & relieve 1 of $g$ & $(a_1a_2,\; a_1a_2g,\; a_1a_2g)$ & $(0,1,1)$ & 2 \\
2 & relieve 2 of $g$ \; [owner$(g) = 3$] & $(a_1a_2,\; a_1a_2,\; a_1a_2g)$ & $(0,0,1)$ & 1 \\
3 & relieve 3 of $a_1$ & $(a_1a_2,\; a_1a_2,\; a_2g)$ & $(0,0,0)$ & 0 \\
4 & relieve 1 of $a_2$ & $(a_1,\; a_1a_2,\; a_2g)$ & $(0,1,1)$ & 2 \\
5 & relieve 2 of $a_1$ \; [owner$(a_1) = 1$] & $(a_1,\; a_2,\; a_2g)$ & $(0,0,1)$ & 1 \\
6 & relieve 3 of $a_2$ \; [owner$(a_2) = 2$] & $(a_1,\; a_2,\; g)$ & $(0,0,0)$ & 0 \\
\bottomrule
\end{tabular}
\end{center}

The terminal allocation is $(\set{a_1},\set{a_2},\set{g})$, the known
zero-subsidy solution. This establishes that the peel template is not killed by
the existing obstruction; it does not establish that the template always works,
which is \Cref{sec:deadends}. Note also how it meets the requirement of
Remark~\ref{rem:deadend} that a correct proof be free to re-open allocated
bundles: \emph{peeling never commits a chore's owner until the last peel of that
type, so ownership is a deferred decision and there is nothing to re-open.}

% --------------------------------------------------------------------------
\subsection{The new obstruction: peel dead ends}
\label{sec:deadends}

The natural conjecture is local.

\begin{conjecture}[Local extendability --- \textbf{refuted}]
\label{conj:h1}
From every legal non-terminal state, some legal move exists.
\end{conjecture}

\begin{proposition}
\label{prop:deadends}
Conjecture~\ref{conj:h1} is false. There are legal, non-terminal peel states
from which no sequence of peels and permutations reaches a terminal state
without violating the invariant.
\end{proposition}

\begin{proof}
A complete backward fixed-point computation over the entire state space of the
witness, with both move types and no heuristics, is carried out by
\texttt{peel3.py}; it reports the following.

\begin{center}
\begin{tabular}{@{}lr@{}}
\toprule
quantity & value \\
\midrule
states in total ($7^3$) & 343 \\
legal states & 175 \\
terminal (partition) states that are legal & 6 of 27 \\
legal states from which a terminal is reachable within the invariant & 166 \\
\textbf{legal dead ends} & \textbf{9} \\
root is good & yes \\
\bottomrule
\end{tabular}
\end{center}

The smallest dead end is $\workprofile = (\set{a_1,a_2},\set{g},\set{g})$, which
is legal with $\pathw{} = (1,0,0)$. The only peelable type is $g$, and both
continuations fail: relieving agent~2 or agent~3 of $g$ creates a weight-$2$
path out of agent~1. Permutations do not help --- the profile
$(\set{g},\set{a_1,a_2},\set{g})$ is legal but its two continuations fail
identically.
\end{proof}

\paragraph{What the dead ends are, in the language of~\cite{BKNS22}.}
This is the sharpest statement available. At
$\workprofile = (\set{a_1,a_2},\set{g},\set{g})$ the remaining unallocated copy
is a copy of $b_g$, and it has three possible recipients. Giving it to agent~2
or~3 is a genuine peel and violates the invariant. Giving it to agent~1 --- who
already holds a copy of $b_g$ --- creates a \emph{duplicate}, which is perfectly
legal in $\hat{I}$ and, by Lemma~\ref{lem:dupfree}, changes nothing whatever in
the graph. So the algorithm of~\cite{BKNS22} never gets stuck on $\hat{I}$: it
gets stuck only in the sense that its one legal move is to spend a copy on a
duplicate, and spending a copy on a duplicate is exactly the loss of coverage.
If agent~1 takes the second $b_g$ then $g$ ends with owner-candidate set
$\set{2,3}$ --- two owners, not a partition. Corollary~\ref{cor:coverage} is
therefore not a formal restatement; the gap it names is realised.

\begin{remark}[The balance signal --- \textbf{observation only}]
\label{rem:balance}
All nine dead ends of the witness have the same shape: \emph{one agent is
already committed to a two-element terminal bundle}. They are exactly the three
choices of which pair times the three choices of which agent. Correspondingly
all six reachable terminals are the six perfect matchings, so on this instance
the only legal partitions are the matchings. This is the same signal as
Theorem~\ref{thm:sizeshift} --- chores are goods plus cardinality balancing ---
and as the balancedness guarantee of~\cite{BDNSV20}, arriving from a third
direction, and it suggests the missing ingredient is a \emph{balance}
discipline rather than a subsidy-based one. Consistent with that, restricting
recipients to $\argmax_i \pathw{}(i)$, the direct mirror of $\Mset(\subsidy)$,
shrinks the reachable state space but does \emph{not} avoid the dead ends: a
subsidy-based rule alone is not enough. This is a single instance and should not
be generalised without more data.
\end{remark}

% --------------------------------------------------------------------------
\subsection{The reduced open problem}
\label{sec:h1prime}

\begin{conjecture}[Reachability --- \textbf{open}]
\label{conj:h1prime}
From the root, some sequence of peels and permutations reaches a terminal state
with the invariant holding at every intermediate state.
\end{conjecture}

Conjecture~\ref{conj:h1prime} implies Conjecture~\ref{conj:main}. The root graph
is identically zero, so the invariant holds there; if a legal schedule reaches a
terminal state $\workprofile^{\mathrm{end}}$ then the induced allocation
$\alloc$ has $\envygraph{\alloc} = \envygraph{\workprofile^{\mathrm{end}}}$,
with no positive cycle and $\pathw{\alloc}(i) \le 1$ for all $i$;
Observation~\ref{obs:integrality} then forces $\optsubsidy \in \set{0,1}^n$, and
the endpoint of a maximum-weight path has $\pathw{} = 0$, giving a total of at
most $n-1$.

\begin{remark}[The converse is open, and it matters for what the experiments test]
\label{rem:converse}
Conjecture~\ref{conj:h1prime} implies Conjecture~\ref{conj:main}, but the
implication has not been established in the other direction and should not be
assumed. Conjecture~\ref{conj:main} asserts that a good terminal state
\emph{exists}; Conjecture~\ref{conj:h1prime} asserts that one is \emph{reachable
along a path every state of which is legal}. These can come apart: a good
terminal state may exist while every schedule leading to it passes through an
illegal intermediate state, and Proposition~\ref{prop:deadends} shows that
illegal-looking intermediate states are not hypothetical. The practical
consequence is that an instance with a \emph{bad root} would refute
Conjecture~\ref{conj:h1prime}, and hence kill this approach, but would say
nothing directly about Conjecture~\ref{conj:main}, which would have to be
attacked by exhaustive search over allocations as in \Cref{sec:experiments}.
Establishing the converse --- that a good terminal state, when one exists, is
always reachable legally --- would make the peel frame a complete
reformulation rather than a sufficient one, and is worth attempting on its own.
\end{remark}

It is not automatic. Conjecture~\ref{conj:h1prime} is a reachability statement
about a state space with genuine dead ends, so a proof must supply a \emph{rule}
for choosing the next peel that provably never enters the bad region, or else
avoid the state space altogether by a potential or exchange argument. Two
resources are available that~\cite{BKNS22} never had:
\begin{enumerate}
  \item \textbf{Scheduling freedom.} Many types are unfinished at once, and we
    need only \emph{some} type to admit a legal peel, not a prescribed one.
  \item \textbf{$\abs{S_j} \ge 2$.} A peelable type always has at least two
    candidate recipients, so the forbidden recipient set is never all of
    $\agents$.
\end{enumerate}

Call a legal non-terminal state a \emph{deadlock} if for every unfinished type
$j$ and every permutation, no recipient in $S_j$ preserves the invariant.
Proposition~\ref{prop:deadends} shows deadlocks exist;
Conjecture~\ref{conj:h1prime} asks only that the \emph{root} avoid them. So the
target theorem has the shape \emph{there is a peel rule under which no deadlock
is ever entered}, and Remark~\ref{rem:balance} is the first candidate for what
that rule must enforce.

% --------------------------------------------------------------------------
\subsection{The type-ordered restriction}
\label{sec:typeorder}

One restriction of the peel process survives every test so far and shrinks the
search substantially. Present the chores in \emph{blocks}: decide the owner of
$a_1$ completely, then $a_2$, and so on. Coverage is automatic in the peel frame
--- a peel deletes $j$ from $\workprofile_x$, so the same pair $(x,j)$ cannot be
peeled twice --- and after $r$ complete blocks the state is
$\workprofile_i = \alloc^{(r)}_i \cup R$ with $R$ the undecided chores, which is
literally the conditioned-remainder principle of Remark~\ref{rem:conditioned}.
A state is then just a triple (partial allocation, current chore, set of agents
already relieved of it), which is a far smaller space than the general peel
process.

\begin{conjecture}[Type-ordered reachability --- \textbf{open}]
\label{conj:h1pp}
From the root, some chore order $\pi$, together with some choice of owner and
within-block sequence for each chore, reaches a terminal state with the
invariant holding at every intermediate state.
\end{conjecture}

Conjecture~\ref{conj:h1pp} implies Conjecture~\ref{conj:h1prime} implies
Conjecture~\ref{conj:main}, and it is a strictly smaller search problem than
either. The evidence, from \texttt{typeorder.py} and \texttt{stress.py}, is that
the restriction costs nothing: across the witness of
Theorem~\ref{thm:obstruction} and $890$ random dichotomous instances at
$n \in \set{3,4}$, $m \le 4$, \emph{every} instance that admits a legal schedule
at all admits a type-ordered one. The chore \emph{order} is not free --- in
$8/120$ instances at $n = m = 3$, $16/60$ at $n=3,m=4$ and $8/60$ at $n=4,m=3$
only some orders work --- but in every instance tested some order does. Note
that the schedule exhibited in \Cref{sec:dissolve} is \emph{not} type-ordered,
since it interleaves $a_1$ and $a_2$; the table says a type-ordered one exists
for that instance anyway, which was not obvious.

What does not survive is the natural recipient rule to pair with it.

\begin{proposition}[$\argmax$ recipient selection is insufficient --- \textbf{refuted}]
\label{prop:f4}
Giving each copy to an agent of maximum current subsidy does not suffice, even
under the type-ordering.
\end{proposition}

\begin{proof}
Take $n = 3$, $m = 2$ with the costs
\[
  \cost_1(\emptyset) = 0,\; \cost_1(\set{a_1}) = \cost_1(\set{a_2}) = 0,\;
  \cost_1(\items) = 1; \qquad
  \cost_2(S) = \min(\abs{S},1); \qquad \cost_3 = \cost_1 .
\]
All three are monotone with marginals in $\set{0,1}$. Exhaustive search over
both chore orders, both roles and all within-block sequences finds no legal
type-ordered schedule in which every recipient lies in
$\argmax_i \pathw{\workprofile}(i)$. Dropping only the $\argmax$ restriction,
the schedule that relieves agent $1$ of $a_1$, then agent $3$ of $a_1$, then
agent $1$ of $a_2$, then agent $3$ of $a_2$ is legal and terminates at
$\alloc = (\emptyset, \set{a_1,a_2}, \emptyset)$ with $\pathw{} = (0,1,0)$.
\end{proof}

This is the signal of Remark~\ref{rem:balance} again --- a subsidy-based
recipient rule alone is not enough --- now confirmed \emph{inside} the
type-ordered restriction, which is the sharper statement. Combined with
Theorem~\ref{thm:f2}, which says the recipient in the hard case has zero
marginal, the rule provably cannot be read off the subsidy vector: it must see
the residual workload, which is exactly what the conditioned-remainder arcs
already encode.

% --------------------------------------------------------------------------
\subsection{A subclass that may fall first --- conjectural}
\label{sec:mrf}

On type space the dual of a \emph{supermodular} dichotomous cost is submodular
with binary marginals, that is, a matroid rank function; lifting through $\tau$
to the copies is the parallel extension of that matroid, still a matroid rank
function. Additive costs dualise to partition-matroid ranks. If that is right
--- and the parallel-extension step in particular deserves a full proof before
being relied on --- then supermodular dichotomous chores are exactly
matroid-rank goods coverage problems, and the witness of
Theorem~\ref{thm:obstruction} lies entirely inside this class, since $\cost_1$
dualises to the rank of $U_{2,3}$, parallel extended.

Matroid rank is where the machinery of the surrounding literature is strongest,
through the exchange properties used in~\cite{BSY23,BEF21}. That makes
Conjecture~\ref{conj:main} restricted to supermodular dichotomous costs, attacked
by basis exchange in the parallel-copy matroid, the natural intermediate target
--- sitting between Theorem~\ref{thm:binadd}, which is closed, and the full
conjecture.
